\documentclass[11pt]{article}
\usepackage[margin=1in]{geometry}
\usepackage{amsmath,amssymb,amsthm}
\usepackage{mathtools}
\usepackage{hyperref}

\newcommand{\eps}{\varepsilon}
\newcommand{\Ax}{A_x}
\newcommand{\Jvec}{\vec J}
\newcommand{\Jreg}{\vec J\,'}

\title{Numerical Evaluation of the Canonical System up to Weight 6\\
\large Regularization at $x\to0$ and Adaptation of Czakon--Tancredi (arXiv:2606.30354)}
\author{}
\date{}

\begin{document}
\maketitle

\section{Strategy of Czakon--Tancredi, and what simplifies here}

Ref.~\cite{CzakonTancredi} solves $d\Jvec = \eps\,A(\vec x)\,\Jvec$ numerically by integrating
along a path, adjoining to the canonical system a set of auxiliary \emph{rational} differential
equations for every non-rational building block (square roots, elliptic periods, \ldots) that
appears in $A$. This avoids ever having to implement the transcendental functions themselves: the
enlarged system is purely rational and can be integrated with a standard adaptive solver.

In the present $371\times371$ system, the connection one-form has, along the line $y\equiv y_0=3/15$,
only \emph{one} algebraic building block, $r(x,y_0)=\sqrt{x\,y_0\,(1-x-y_0)}$. Unlike an elliptic
period, $r$ is an ordinary algebraic function with a closed form on the real line for
$0<x<1-y_0$ (no differential equation needs to be adjoined for it: it is simply evaluated). The
only genuine complication is that $x=0$ is a regular singular (Fuchsian) point of the system, so the
paper's ``integrate directly'' strategy must be preceded by the regularization worked out in
Sections~\ref{sec:setup}--\ref{sec:recover} below.

\section{Setup}
\label{sec:setup}

The canonical system is $D\Jvec = \eps\,\tilde A\,\Jvec$, where $\tilde A(x,y)$ (\texttt{Atilde.m}) is a
matrix of logarithms of alphabet letters, so that the genuine connection one-form is obtained by
differentiating entrywise,
\begin{equation}
A_x^{\rm full}(x,y) \equiv \frac{\partial \tilde A}{\partial x} = \texttt{dx.m}\,, \qquad
A_y^{\rm full}(x,y) \equiv \frac{\partial \tilde A}{\partial y} = \texttt{dy.m}\,.
\end{equation}
Restricting to the line $y \equiv y_0 = 3/15$ and parametrizing $x=x(t)$, the system collapses to a
genuine ODE in $t$,
\begin{equation}
\frac{d\Jvec}{dt} = \eps\,\dot x(t)\,A_x^{\rm full}\big(x(t),y_0\big)\,\Jvec(t)\,,
\end{equation}
since $\dot y = 0$ along this path.

The point $x=0$ is a Fuchsian singularity of $A_x^{\rm full}$. Splitting off the pole exactly,
\begin{equation}
A_x^{\rm full}(x,y_0) = \frac{\Ax}{x} + B(x,y_0)\,, \qquad
\Ax \equiv \operatorname*{Res}_{x=0} A_x^{\rm full}(x,y_0) = \texttt{toexp1.m}\,,
\end{equation}
where $B(x,y_0) \equiv A_x^{\rm full}(x,y_0) - \Ax/x$ is regular (finite) as $x\to0$. Here $\Ax$ is a
constant $371\times371$ matrix, verified to coincide exactly with \texttt{toexp1.m} and with the
residue of \texttt{dx.m} at $x=0$.

\section{Regularization at $x\to0$}

Define
\begin{equation}
M(x,\eps) \equiv \exp\!\big(\eps \Ax \log x\big) = x^{\eps \Ax}\,,
\end{equation}
which is exactly \texttt{exponentiated1.m} with $\delta \to x$, and the \emph{regularized vector}
\begin{equation}
\Jreg(x) \equiv M(x,\eps)^{-1}\Jvec(x)\,, \qquad\text{equivalently}\qquad \Jvec(x) = M(x,\eps)\,\Jreg(x)\,.
\end{equation}
Since $M(x,\eps)$ is built purely from $\Ax$, it commutes with $\Ax$ and
\begin{equation}
\frac{dM}{dx} = \frac{\eps}{x}\,\Ax\,M(x,\eps) = \frac{\eps}{x}\,M(x,\eps)\,\Ax\,.
\end{equation}
Differentiating $\Jvec = M\Jreg$ and imposing the original equation,
\begin{equation}
\frac{\eps}{x}\Ax M \Jreg + M\frac{d\Jreg}{dx}
= \eps\left(\frac{\Ax}{x} + B(x,y_0)\right) M \Jreg\,.
\end{equation}
The $\Ax/x$ terms cancel \emph{identically} on both sides — this cancellation is exactly what fixes
the sign of the exponent in $M(x,\eps)$ to be $+\eps \Ax$, matching the convention already used to build
\texttt{exponentiated1.m}; no sign flip is required. What remains is
\begin{equation}
M(x,\eps)\,\frac{d\Jreg}{dx} = \eps\, B(x,y_0)\, M(x,\eps)\, \Jreg(x)\,,
\end{equation}
i.e.\ the \textbf{exact, pole-free equation}
\begin{equation}
\boxed{\ \dfrac{d\Jreg}{dx} = \eps\, M(x,\eps)^{-1} B(x,y_0)\, M(x,\eps)\, \Jreg(x)\ }\,, \qquad
\frac{d\Jreg}{dt} = \dot x(t) \times (\text{same, at } x=x(t))\,,
\label{eq:regODE}
\end{equation}
with
\begin{equation}
B(x,y_0) = \texttt{dx.m}(x,y_0) - \frac{\texttt{toexp1.m}}{x}
\end{equation}
manifestly finite (no pole) at $x=0$. The conjugated coefficient $M(x,\eps)^{-1}B(x,y_0)M(x,\eps)$ remains
bounded as $x\to 0$, order by order in $\eps$: any $\log x$ dependence it carries integrates to terms
like $x\log x \to 0$, never to a divergence. Hence this equation may be integrated numerically
starting exactly at $x=0$, using the known boundary value $\Jreg(0,y_0)$.

\section{Recovering the physical vector}
\label{sec:recover}

At the target point $x_1$,
\begin{equation}
\Jvec(x_1,y_0) = M(x_1,\eps)\,\Jreg(x_1,y_0)
= \texttt{exponentiated1.m}\big|_{\delta = x_1} \cdot \Jreg(x_1,y_0)\,.
\end{equation}

\paragraph{Practical note on $M(x,\eps)^{-1}$.} No new computation is required: since
$M(x,\eps)^{-1} = \exp(-\eps \Ax \log x)$ is built from the \emph{same} spectral projectors and nilpotents
already saved in \texttt{toJD/Npowers\_$\langle\lambda\rangle$\_toexp1.m}, it is obtained by flipping
the sign of $\eps$ (equivalently $\lambda \to -\lambda$ in $\delta^{\eps\lambda}$ and
$\eps\log\delta \to -\eps\log\delta$ in the nilpotent series).

\section{Weight-by-weight (order-by-order in $\eps$) reduction}

Rather than solving \eqref{eq:regODE} at a fixed numerical $\eps$, we extract the Taylor
coefficients in $\eps$ directly -- this is both what is physically wanted (the finite remainder is
built order by order in $\eps$) and what makes the numerics cheap, since $M(x,\eps)$ has an
\emph{exactly known}, finite eps-expansion built from already-computed exact matrices.

\subsection{Expansion of $M(x,\eps)$}

Since $A_x=\texttt{toexp1.m}$ has nilpotency order $\le 2$ (only at $\lambda=-3$; order $1$ at
$\lambda=1$; semisimple elsewhere), for each eigenvalue $\lambda$
\begin{equation}
x^{\eps\lambda}\Big(P_\lambda + \eps\log x\, N_\lambda + \tfrac{1}{2}(\eps\log x)^2 N_\lambda^2\Big)
\end{equation}
is a product of the ordinary scalar exponential series $x^{\eps\lambda}=\sum_n \frac{(\lambda\log
x)^n}{n!}\eps^n$ with a \emph{polynomial} of degree $\le2$ in $\eps$. Convolving the two series and
summing over $\lambda$ gives, for any truncation order $W$,
\begin{equation}
M(x,\eps) = \sum_{j=0}^{W-1} \eps^j M_j(x) + O(\eps^W)\,, \qquad
M_j(x) = \sum_\lambda\ \sum_{k=0}^{\min(2,j)} \frac{(\lambda\log x)^{j-k}}{(j-k)!}\cdot
\frac{(\log x)^k}{k!}\,N_\lambda^k\,.
\label{eq:Mj}
\end{equation}
Every $M_j(x)$ is an explicit \emph{linear combination of the fixed, exactly-known matrices}
$\{P_\lambda, N_\lambda, N_\lambda^2\}$ (already saved in \texttt{toJD/Npowers\_$\langle\lambda\rangle$\_toexp1.m})
with scalar coefficients that are simple powers of $\log x$. No matrix multiplication is needed to
build $M_j(x)$ -- only scaled matrix additions. In fact, since $(\log x)^{j-k}(\log x)^k=(\log x)^j$
regardless of $k$, this simplifies further to $M_j(x) = (\log x)^j\, Q_j$ with
\begin{equation}
Q_j \equiv \sum_\lambda\ \sum_{k=0}^{\min(2,j)} \frac{\lambda^{j-k}}{(j-k)!\,k!}\,N_\lambda^k
\end{equation}
a fixed, \emph{$x$-independent} exact rational matrix, computable once ($Q_0=\sum_\lambda P_\lambda
= \mathbb 1$, $Q_1 = \sum_\lambda(\lambda P_\lambda + N_\lambda)=\Ax$, verified explicitly).

\subsection{Expansion of $M(x,\eps)^{-1}$}

Since $M(x,\eps)^{-1}=\exp(-\eps A_x\log x)=M(x,-\eps)$,
\begin{equation}
M(x,\eps)^{-1} = \sum_{j=0}^{W-1}\eps^j\,(-1)^j M_j(x) + O(\eps^W)\,,
\end{equation}
i.e.\ the inverse's coefficients are obtained from \eqref{eq:Mj} by a sign flip on odd $j$ -- no
independent computation is required.

\subsection{Expansion of the conjugated matrix and the hierarchy for $\Jreg$}

Write $C(x,\eps)=M(x,\eps)^{-1}B(x,y_0)M(x,\eps)=\sum_{j\ge0}\eps^j C_j(x)$, with
\begin{equation}
C_j(x) = \sum_{p+q=j} (-1)^p M_p(x)\,B(x,y_0)\,M_q(x)
= (\log x)^j \sum_{p=0}^{j} (-1)^p Q_p\,B(x,y_0)\,Q_{j-p}\,,
\end{equation}
using $M_p(x)M_q(x)\propto(\log x)^{p+q}=(\log x)^j$ from the previous subsection. Expanding
$\Jreg(x,\eps)=\sum_{k\ge0}\eps^k \Jreg^{(k)}(x)$ and matching powers of $\eps$ in
\eqref{eq:regODE} gives the \textbf{exact} recursive hierarchy
\begin{equation}
\boxed{\ \frac{d\Jreg^{(0)}}{dx}=0\,,\qquad
\frac{d\Jreg^{(m)}}{dx} = \sum_{k=0}^{m-1} C_{m-1-k}(x)\,\Jreg^{(k)}(x)\,,\quad m=1,\dots,W-1\ }
\label{eq:hierarchy}
\end{equation}
Crucially, the right-hand side at order $m$ involves \emph{only} $\Jreg^{(0)},\ldots,\Jreg^{(m-1)}$,
never $\Jreg^{(m)}$ itself. The system is therefore \emph{block lower-triangular}: $\Jreg^{(0)}$ is
the constant weight-0 boundary vector (rational, valid everywhere -- consistent with
\texttt{sol0.m}), and each subsequent order is obtained by integrating a source built entirely from
lower orders. This is precisely the iterated-integral structure of the canonical system, made
explicit rather than left implicit in a black-box $\eps$-integration.

\subsection{Practical evaluation of $B(x,y_0)$}

$B(x,y_0)=\texttt{dx.m}(x,y_0)-A_x/x$ is $\eps$-independent. Along the fixed line $y=y_0$, its only
non-rational content is the algebraic letter $r(x,y_0)=\sqrt{x\,y_0(1-x-y_0)}$ and the derived
combinations $r_1=r(x-1)/y_0$, $r_2=r(x+y_0)/y_0$ appearing in \texttt{dx.m}; substituting these
closed forms turns $B(x,y_0)$ into an explicit (compilable) function of $x$ alone. No auxiliary ODE
is needed for $r$ on this path, since it is not a period of a higher-genus curve but an elementary
algebraic function.

\section{Numerical integration}

The state vector is the stack $\big(\Jreg^{(1)},\ldots,\Jreg^{(W-1)}\big)\in\mathbb R^{371(W-1)}$
($\Jreg^{(0)}$ is constant and known, so it is not part of the dynamical state). Each block evolves
via \eqref{eq:hierarchy}, with the $C_j(x)$ evaluated on the fly at each requested $x$ (via
$Q_j$, cheap scaled sums, plus the compiled $B(x,y_0)$ and $371\times371$ matrix products
per $C_j$). This is integrated with \texttt{NDSolve} (adaptive Runge--Kutta by default), starting
from a small $x=x_{\rm start}\approx0$ (a numerical regulator for $\log x$; $\Jreg$ is analytic at
$x=0$ so $\Jreg(x_{\rm start})\simeq\Jreg(0)$ to excellent approximation) up to the target point
$x_1$. Finally, the physical vector at $x_1$ is recovered by re-multiplying and
re-expanding in $\eps$ up to $O(\eps^{W-1})$:
\begin{equation}
\Jvec(x_1,y_0) = M(x_1,\eps)\,\Jreg(x_1,\eps)\Big|_{\eps\text{-expanded to order }W-1}
= \sum_{m=0}^{W-1}\eps^m \sum_{k=0}^m M_{m-k}(x_1)\,\Jreg^{(k)}(x_1)\,.
\end{equation}
This is implemented in \texttt{integrate\_canonical.wl}, parameterized by the target point $x_1$ and
the maximal weight $W-1\le6$.

\begin{thebibliography}{9}
\bibitem{CzakonTancredi} M.~Czakon and L.~Tancredi,
\emph{Solution of Canonical Differential Equations for Integrals on Arbitrary Geometries},
arXiv:2606.30354 [hep-ph].
\end{thebibliography}

\end{document}
